\section*{Capítulo \ref{ch:g6:describing-motion} --- Descrever o movimento: trajetória e velocidade}

\begin{solution}{exo:g6:describing-motion:1}
Retilínea (descendo a vidraça); circular; curva (um arco); curva (um
S serpenteante).
\end{solution}

\begin{solution}{exo:g6:describing-motion:2}
O eixo: uma linha reta, deslizando paralela à rua. A válvula: uma
curva de saltinhos graciosos --- subindo, arqueando e tocando o chão
a cada volta da roda.
\end{solution}

\begin{solution}{exo:g6:describing-motion:3}
Uniforme; variado (acelerando); uniforme; variado (freando).
\end{solution}

\begin{solution}{exo:g6:describing-motion:4}
Vãos iguais: \omterm{def:g6:describing-motion:uniform}{movimento uniforme}. Vãos crescendo ($1, 2, 4, 7$):
acelerando --- mais distância em cada batida igual.
\end{solution}

\begin{solution}{exo:g6:describing-motion:5}
Uma câmera que se desloca acrescenta o movimento dela ao registro ---
os pontos então misturam a viagem do objeto com a da câmera. O
movimento é sempre movimento em relação a alguma coisa, e o filme
precisa fixar essa ``alguma coisa'': é a observação sobre pontos de
vista.
\end{solution}

\begin{solution}{exo:g6:describing-motion:6}
Tabela: $30$ min $\to$ \qty{12}{km}, então $60$ min $\to$
\qty{24}{km}: a \omterm{def:g5:speed-distance-time:speed}{velocidade} da balsa é de $24$ quilômetros por hora.
\end{solution}

\begin{solution}{exo:g6:describing-motion:7}
$60$ min $\to$ \qty{14}{km}; $30$ min $\to$ \qty{7}{km}; $15$ min
$\to$ \qty{3.5}{km}.
\end{solution}

\begin{solution}{exo:g6:describing-motion:8}
Um passeante a $12$ quilômetros por hora: rejeitado --- três vezes o
ritmo de passeio. O ônibus: $10$ km em $30$ min dá $20$ quilômetros
por hora --- exatamente o ritmo de ônibus urbano, aceito. O trem:
$75$ km em $15$ min dá $300$ quilômetros por hora --- topo do quadro,
mas plausível para um trem de alta \omterm{def:g5:speed-distance-time:speed}{velocidade}: aceito.
\end{solution}

\begin{solution}{exo:g6:describing-motion:9}
Para o passageiro à frente: nenhum movimento --- nenhuma \omterm{def:g6:describing-motion:trajectory}{trajetória},
ritmo zero. Para a vaca: uma longa \omterm{def:g6:describing-motion:trajectory}{trajetória retilínea} no ritmo
completo do trem. As duas descrições são honestas: cada uma descreve
o passageiro que cochila em relação ao próprio ponto de vista, e
movimento é sempre movimento em relação a alguma coisa.
\end{solution}

\begin{solution}{exo:g6:describing-motion:10}
Tempo de deslocamento: $18 - 6 = 12$ minutos para \qty{9}{km}.
Tabela: $12$ min $\to$ $9$ km, então $60$ min $\to$ $45$ km ---
exatamente $45$ quilômetros por hora: nem acima nem abaixo, mas
certinho.
\end{solution}

\begin{solution}{exo:g6:describing-motion:11}
Filme a busca da bolinha com a câmera parada; pause a cada meio
segundo; marque as posições do cachorro em relação às estacas da
cerca; meça os vãos. Vãos iguais o tempo todo: ritmo constante. Vãos
que se abrem e depois encolhem (ou que alternam largos e estreitos):
dispara e desliza. Os pontos entregam o veredito, não importa o que a
família tenha gritado.
\end{solution}

\begin{solution}{exo:g6:describing-motion:12}
Da sala: uma curva --- a bola arqueia para a frente e para baixo, da
beirada da mesa até o chão. Do carrinho que rola no ritmo do avanço
da bola: a bola não ganha nem perde terreno --- ela parece cair
\emph{em linha reta para baixo}. A mesma queda, dois pontos de vista,
e a vista mais simples pertence ao observador em movimento: uma
semente de uma grande ideia.
\end{solution}

\begin{solution}{pb:g6:describing-motion:1}
\textbf{1.} Vãos iguais de \qty{50}{m}: \omterm{def:g6:describing-motion:uniform}{movimento uniforme}.
\textbf{2.} $5 \times 50 = \qty{250}{m}$.
\textbf{3.} O bonde arranca do repouso: vãos curtos que vão
crescendo --- acelerando nas três primeiras batidas --- e depois se
acomodam em vãos constantes de \qty{50}{m}: cruzeiro.
\textbf{4.} Nos dois últimos intervalos, de \qty{50}{m} cada.
\textbf{5.} $10$ s $\to$ \qty{50}{m}, então $60$ s $\to$
\qty{300}{m} por minuto.
\textbf{6.} $60$ min $\to$ $60 \times 300 = \qty{18000}{m}$:
\omterm{def:g5:speed-distance-time:speed}{velocidade} de cruzeiro de $18$ quilômetros por hora.
\textbf{7.} Tabela: \qty{18}{km} por $60$ min, então \qty{12}{km} em
$40$ minutos.
\textbf{8.} $50 - 40 = 10$ minutos perdidos com paradas e trechos
lentos.
\textbf{9.} $50$ min $\to$ \qty{12}{km}, então $60$ min $\to$
$12 \times 60 \div 50 = \qty{14.4}{km}$: a viagem inteira corre a
cerca de $14$ quilômetros por hora --- os $18$ do cartaz são honestos
só para os trechos de cruzeiro, e não para ``atravessar a cidade''.
\textbf{10.} Uma \omterm{def:g5:speed-distance-time:speed}{velocidade} de viagem inteira descreve o percurso
\emph{como se} ele fosse uniforme: o único ritmo constante que
cobriria os mesmos \qty{12}{km} nos mesmos $50$ minutos. Velocímetro
nenhum a mostra e, ainda assim, ela é o número justo para o
planejamento, porque já contém cada parada e cada trecho lento.
\textbf{11.} $50 - 4 = 46$ minutos;
$12 \times 60 \div 46 \approx \qty{15.7}{km}$ por hora --- digamos
uns $16$.
\textbf{12.} Por exemplo: ``Os passageiros devem se planejar pela
\omterm{def:g5:speed-distance-time:speed}{velocidade} de viagem inteira --- cerca de $14$ quilômetros por hora
hoje --- porque ela é o único número que inclui as esperas; os $18$
de cruzeiro bajulam a linha e ignoram todas as paradas.''
\end{solution}
